\chapter{Introduction}

\section{Purpose of this course}
To calculate the responses, including internal forces (axial, shear, moment) distribution and displacements (deflections), of framed structures.

\section{Five elements of structural analysis}
\begin{enumerate}
  \item \textbf{Basic mechanics:}
    \begin{enumerate}[label=(\alph*)]
      \item stress--strain relationships
      \item equilibrium equations
      \item compatibility equations
    \end{enumerate}
  \item \textbf{Finite element mechanics}, or member force--deformation relationships
    \begin{fig}
      % Simply supported beam with end moments M1(alpha1), M2(alpha2)
\begin{tikzpicture}[scale=1.1,>=Latex]
  \def\L{6}
  % beam
  \draw[line width=2.2pt,black!70] (0,0) -- (\L,0);
  \node[above] at (\L/2,0.05) {$EI$};
  % pin support (left)
  \draw[fill=black!60] (0,0) -- (-0.25,-0.4) -- (0.25,-0.4) -- cycle;
  \fill[white,draw=black] (0,0) circle (0.06);
  \fill[black!25] (-0.4,-0.4) rectangle (0.4,-0.55);
  % roller support (right)
  \draw[fill=black!60] (\L,0) -- (\L-0.25,-0.4) -- (\L+0.25,-0.4) -- cycle;
  \fill[white,draw=black] (\L,0) circle (0.06);
  \draw[fill=white] (\L-0.15,-0.5) circle (0.09);
  \draw[fill=white] (\L+0.15,-0.5) circle (0.09);
  \fill[black!25] (\L-0.4,-0.6) rectangle (\L+0.4,-0.75);
  % end moments
  \draw[->,thick] (-0.05,0.55) arc (150:-30:0.35);
  \node[above] at (0.2,0.75) {$M_1(\alpha_1)$};
  \draw[->,thick] (\L+0.05,0.55) arc (30:210:0.35);
  \node[above] at (\L-0.2,0.75) {$M_2(\alpha_2)$};
  % length dimension
  \draw[<->] (0,-1.1) -- (\L,-1.1) node[midway,fill=white,inner sep=1pt] {$L$};
  \draw (0,-0.85) -- (0,-1.3);
  \draw (\L,-0.85) -- (\L,-1.3);
\end{tikzpicture}

    \end{fig}
    \begin{itemize}
      \item \emph{Stiffness methods}: forces in terms of displacements through stiffness coefficients
        \begin{equation}
          \Bmat{M_1\\ M_2} = \frac{2EI}{L}\bmat{2 & 1\\ 1 & 2}\Bmat{\alpha_1\\ \alpha_2}
        \end{equation}
      \item \emph{Flexibility methods}: deformations in terms of forces through flexibility coefficients
        \begin{equation}
          \Bmat{\alpha_1\\ \alpha_2} = \frac{L}{6EI}\bmat{2 & -1\\ -1 & 2}\Bmat{M_1\\ M_2}
        \end{equation}
    \end{itemize}
  \item \textbf{Formulation of structural (global) equations}\\
    Meaning of structural equations:
    \begin{itemize}
      \item Stiffness methods: equilibrium equations
      \item Flexibility methods: compatibility equations (consistent deformation)
    \end{itemize}
  \item \textbf{Equation solving}\\
    Gauss elimination, decomposition methods, etc.
  \item \textbf{Solution interpretation}
    \begin{itemize}
      \item Member force diagrams such as axial, shear and moment diagrams
      \item Deformed shapes
    \end{itemize}
\end{enumerate}

\section{Types of analysis}
(See also Sec.~8.1 of the textbook.)

\subsection*{Linear analysis}

\subsection*{Nonlinear analysis}

\topicnote{Geometrically nonlinear analysis}{(second-order analysis, $P$-$\Delta$ effects)}
\begin{fig}
  % P-Delta effect on a cantilever column: three stages
\begin{tikzpicture}[scale=1.0,>=Latex,font=\small]
  \def\h{3.2}
  % ---------- (a) undeformed ----------
  \begin{scope}
    \fill[black!45] (-0.55,-0.35) rectangle (0.55,0);
    \draw[fill=white] (-0.09,0) rectangle (0.09,\h);
    \draw[->,thick] (0,\h+0.75) -- (0,\h+0.08) node[pos=0,above] {$P$};
    \draw[->,thick] (-1.0,\h) -- (-0.15,\h) node[pos=0,left] {$V$};
    \draw[<->] (-1.3,0) -- (-1.3,\h) node[midway,left] {$h$};
    % reactions
    \draw[->,thick] (1.1,-0.18) -- (0.6,-0.18) node[pos=0,right] {$V$};
    \draw[->,thick] (0,-0.8) -- (0,-0.4) node[pos=0,below] {$P$};
    \draw[->,thick] (-0.75,-0.75) arc (-120:20:0.3);
    \node[left] at (-0.8,-0.6) {$Vh$};
  \end{scope}
  \draw[-{Triangle[width=6pt,length=6pt]},line width=1.6pt] (1.6,\h/2) -- (2.3,\h/2);
  % ---------- (b) first-order deformed ----------
  \begin{scope}[xshift=4.2cm]
    \fill[black!45] (-0.55,-0.35) rectangle (0.55,0);
    \draw[thick] (0,0) .. controls (0,\h*0.55) and (0.15,\h*0.85) .. (0.55,\h);
    \draw[->,thick] (-0.9,\h) -- (0.4,\h) node[pos=0,left] {$V$};
    \draw[->,thick] (1.1,-0.18) -- (0.6,-0.18) node[pos=0,right] {$V$};
    \draw[->,thick] (-0.75,-0.75) arc (-120:20:0.3);
    \node[left] at (-0.8,-0.6) {$Vh$};
  \end{scope}
  % ---------- (c) second-order deformed ----------
  \begin{scope}[xshift=8.8cm]
    \fill[black!45] (-0.55,-0.35) rectangle (0.55,0);
    \draw[densely dashed] (0,0) .. controls (0,\h*0.55) and (0.15,\h*0.85) .. (0.55,\h);
    \draw[thick] (0,0) .. controls (0,\h*0.5) and (0.2,\h*0.85) .. (0.85,\h);
    \draw[->,thick] (0.85,\h+0.75) -- (0.85,\h+0.08) node[pos=0,above] {$P$};
    \draw[->,thick] (-0.9,\h) -- (0.65,\h) node[pos=0,left] {$V$};
    \draw[->,thick] (1.1,-0.18) -- (0.6,-0.18) node[pos=0,right] {$V$};
    \draw[->,thick] (0,-0.8) -- (0,-0.4) node[pos=0,below] {$P$};
    \draw[->,thick] (-0.75,-0.75) arc (-120:20:0.3);
    \node[left] at (-0.8,-0.6) {$Vh+P\Delta$};
    \draw[<->] (0,\h+0.35) -- (0.85,\h+0.35) node[midway,above] {$\Delta$};
  \end{scope}
  % ---------- formulas ----------
  \node[anchor=west,align=left] at (11.0,\h*0.55) {%
    $\Delta \cong \Delta_L \times \dfrac{1}{1-\dfrac{P}{P_E}}$\\[10pt]
    $\Delta_L=\dfrac{Vh^3}{3EI}$\\[8pt]
    $P_E=\dfrac{\pi^2EI}{4h^2}$};
\end{tikzpicture}

\end{fig}

\topic{Materially nonlinear analysis}
\begin{fig}
  % Materially nonlinear analysis: nonlinear elastic vs. plastic
\begin{tikzpicture}[scale=1.0,>=Latex,font=\small]
  % ---------- nonlinear elastic ----------
  \begin{scope}
    \draw[->] (0,0) -- (3.4,0) node[right] {$\varepsilon$};
    \draw[->] (0,0) -- (0,2.6) node[above] {$\sigma$};
    \draw[thick] (0,0) .. controls (0.6,1.7) and (1.6,2.2) .. (3.0,2.3);
    % loading / unloading arrows along the same curve
    \draw[->,thin] (0.35,0.85) .. controls (0.5,1.2) .. (0.85,1.55);
    \draw[<-,thin] (1.3,1.83) .. controls (1.5,1.95) .. (1.9,2.05);
    \node[left] at (-0.2,1.3) {\itshape elastic};
  \end{scope}
  % ---------- plastic ----------
  \begin{scope}[xshift=6.8cm]
    \draw[->] (0,0) -- (3.4,0) node[right] {$\varepsilon$};
    \draw[->] (0,0) -- (0,2.6) node[above] {$\sigma$};
    % hysteresis parallelogram, hatched
    \fill[pattern={Lines[angle=90,distance=3pt]},pattern color=black!70]
      (0,0) -- (0.8,1.9) -- (2.4,1.9) -- (1.6,0) -- cycle;
    \draw[thick] (0,0) -- (0.8,1.9) -- (2.4,1.9) -- (1.6,0) -- cycle;
    \draw[->,thin] (0.35,0.83) -- (0.5,1.19);
    \draw[->,thin] (1.4,1.9) -- (1.7,1.9);
    \draw[->,thin] (2.05,1.07) -- (1.9,0.71);
    \node[left] at (-0.2,1.1) {\itshape plastic};
  \end{scope}
\end{tikzpicture}

\end{fig}

\topicnote{Nonlinear analysis due to change in boundary conditions}{(such as contact problems)}
\begin{fig}
  % Contact problem: a ball pressed onto a rigid surface
\begin{tikzpicture}[scale=1.0,>=Latex]
  % ground
  \fill[pattern={Lines[angle=45,distance=4pt]},pattern color=black!70] (-1.6,-0.35) rectangle (1.6,0);
  \draw[thick] (-1.6,0) -- (1.6,0);
  % ball (shaded)
  \shade[ball color=black!45] (0,1.0) circle (1.0);
  % load
  \draw[->,very thick] (0,3.0) -- (0,2.1);
  \node[right] at (0.1,2.75) {$P(\Delta)$};
\end{tikzpicture}

\end{fig}
Considering a ball subjected to a load $P$, the contact area will increase as the load increases.

\subsection*{Example of geometrically nonlinear analysis}
A simple example, the \emph{two-member truss}, illustrating the geometrically nonlinear analysis is as follows.
\begin{fig}
  % Two-member (von Mises) truss: geometrically nonlinear example
\begin{tikzpicture}[scale=1.0,>=Latex]
  \def\Lx{5.0} \def\h{2.0} \def\dl{0.7}
  \coordinate (A) at (0,0);
  \coordinate (B) at (2*\Lx,0);
  \coordinate (C0) at (\Lx,\h);
  \coordinate (C1) at (\Lx,\h-\dl);
  % undeformed (dotted)
  \draw[densely dotted] (A) -- (C0) -- (B);
  % baseline (dashed)
  \draw[dashed] (A) -- (B);
  % deformed (thick)
  \draw[line width=2.5pt,black!70] (A) -- (C1) -- (B);
  % supports
  \foreach \p in {A,B} {
    \draw[fill=black!60] (\p) -- ++(-0.25,-0.4) -- ++(0.5,0) -- cycle;
    \fill[black!25] ($(\p)+(-0.4,-0.4)$) rectangle ($(\p)+(0.4,-0.55)$);
    \fill[white,draw=black] (\p) circle (0.06);
  }
  \fill[white,draw=black] (C1) circle (0.07);
  % load
  \draw[->,very thick] ($(C0)+(0,1.0)$) -- ($(C1)+(0,0.1)$);
  \node[above] at ($(C0)+(0,1.0)$) {$P$};
  % labels
  \node[above left,inner sep=1pt] at ($(A)!0.55!(C0)$) {$L$};
  \node[above,inner sep=1pt] at ($(A)!0.62!(C1)$) {$L'$};
  % angles
  \draw[->] ($(A)+(1.4,0)$) arc (0:{atan(\h/\Lx)}:1.4);
  \node at ($(A)+(1.65,0.22)$) {$\theta$};
  \draw[->] ($(A)+(2.3,0)$) arc (0:{atan((\h-\dl)/\Lx)}:2.3);
  \node at ($(A)+(2.55,0.22)$) {$\theta'$};
  % dimension lines
  \draw[thin] (C0) -- ($(B)+(1.0,\h)$);
  \draw[thin] (C1) -- ($(B)+(1.0,\h-\dl)$);
  \draw[thin] (B) -- ($(B)+(1.0,0)$);
  \draw[<->] ($(B)+(0.6,\h)$) -- ($(B)+(0.6,\h-\dl)$) node[midway,right] {$\Delta$};
  \draw[<->] ($(B)+(0.9,\h-\dl)$) -- ($(B)+(0.9,0)$) node[midway,right] {$h-\Delta$};
\end{tikzpicture}

\end{fig}

\begin{solution}
For a displacement $\Delta$, we have
\begin{equation}
  L' = \sqrt{(h-\Delta)^2 + (L\cos\theta)^2}.
\end{equation}
Knowing $\eps = \dfrac{L-L'}{L}$ and the stress--strain relationship $\sigma = f(\eps)$, we can first obtain the stress corresponding to this displacement and then the axial force $F = \sigma A$. Finally, the equilibrium equation states that
\begin{equation}
  P = 2F\sin\theta', \qquad \text{where } \sin\theta' = (h-\Delta)/L'.
\end{equation}
It shall be noted that the equilibrium equation is written for the \emph{deformed configuration}, which is the main feature of geometric nonlinear analysis. With different stress--strain relationships, the $P$--$\Delta$ curves are of different shapes, but all the curves pass through $(0,0)$, $(h,0)$ and $(2h,0)$.
\end{solution}
\begin{fig}
  % Load-displacement curve of the two-member truss (snap-through)
\begin{tikzpicture}[>=Latex]
  \begin{axis}[
    width=12cm,height=6cm,
    axis lines=middle,
    xlabel={$\Delta$},ylabel={$P$},
    xlabel style={at={(ticklabel* cs:1)},anchor=west},
    ylabel style={at={(ticklabel* cs:1)},anchor=south},
    xmin=0,xmax=2.6,ymin=-1.4,ymax=1.9,
    xtick={0.1,1,2},xticklabels={$\Delta_i$,$h$,$2h$},
    ytick=\empty,
    clip=false,
    samples=120,domain=0:2.35,
  ]
    % P proportional to (h-D)(D)(2h-D) (shape of the exact curve)
    \addplot[thick,black] {2.6*(1-x)*x*(2-x)};
    % P_lim and P_i
    \pgfmathsetmacro{\plim}{2.6*(1-0.4226)*0.4226*(2-0.4226)}
    \pgfmathsetmacro{\pI}{2.6*(1-0.1)*0.1*(2-0.1)}
    \draw[densely dotted] (axis cs:0,\plim) -- (axis cs:2.2,\plim);
    \draw[densely dotted] (axis cs:0,\pI) -- (axis cs:0.1,\pI) -- (axis cs:0.1,0);
    \node[left] at (axis cs:0,\plim) {$P_{\lim}$};
    \node[left] at (axis cs:0,\pI) {$P_i$};
  \end{axis}
\end{tikzpicture}

\end{fig}

Another famous example related to geometric nonlinear analysis is a horizontal cable fixed at both ends and with a pretension $T$. Assume the length of the cable is $2L$ and the cable is subjected to a vertical load $P$ applied at the midpoint. Assuming $T$ is much larger than $P$, calculate the vertical deflection at the loaded point.
\begin{fig}
  % Pretensioned cable with a concentrated load at midspan
\begin{tikzpicture}[scale=1.0,>=Latex]
  \def\L{4.0} \def\v{0.8}
  \coordinate (A) at (0,0); \coordinate (B) at (2*\L,0); \coordinate (C) at (\L,-\v);
  % undeformed cable (dotted)
  \draw[densely dotted] (A) -- (B);
  % deformed cable (double line)
  \draw[double,double distance=1.5pt] (A) -- (C) -- (B);
  % supports
  \foreach \p in {A,B} {
    \draw[fill=black!60] (\p) -- ++(-0.25,-0.4) -- ++(0.5,0) -- cycle;
    \fill[black!25] ($(\p)+(-0.4,-0.4)$) rectangle ($(\p)+(0.4,-0.55)$);
    \fill[white,draw=black] (\p) circle (0.06);
  }
  \draw[fill=white] (C) circle (0.12); \fill (C) circle (0.05);
  % tension arrows (cut)
  \draw[->,very thick] ($(A)+(-0.3,0)$) -- ++(-1.1,0) node[left] {$T$};
  \draw[->,very thick] ($(B)+(0.3,0)$) -- ++(1.1,0) node[right] {$T$};
  % load
  \draw[->,very thick] (C) -- ++(0,-1.0) node[right] {$P$};
  % labels
  \node[below left] at ($(A)!0.5!(C)$) {$L'$};
  \draw[->] ($(A)+(2.2,0)$) arc (0:{-atan(\v/\L)}:2.2);
  \node[right] at ($(A)+(2.25,-0.2)$) {$\theta$};
  % dimensions
  \draw[<->] (-0.7,0) -- (-0.7,-\v) node[midway,left] {$v$};
  \draw[thin] (-0.45,0) -- (-0.9,0); \draw[thin] (C) -- (-0.9,-\v);
  \draw[<->] (0,-1.4) -- (\L,-1.4) node[midway,fill=white,inner sep=1pt] {$L$};
  \draw[thin] (0,-1.2) -- (0,-1.6); \draw[thin] (\L,-1.2) -- (\L,-1.6);
\end{tikzpicture}

\end{fig}

\begin{solution}
The equilibrium equation for the deformed configuration is
\begin{equation}
  2(T+\Delta F)\sin\theta = P,
\end{equation}
where $\Delta F$ is the increment of the tension. Since $T$ is much larger than $P$, $\Delta F$ must be very small compared with $T$. Also, $\sin\theta \approx v/L$; hence,
\begin{equation}
  2T\frac{v}{L} = P \quad\Longrightarrow\quad v = \frac{PL}{2T}.
\end{equation}
The relation $v \propto \dfrac{P}{T}$ is what we would expect.
\end{solution}

\section{Different levels of analysis}
\begin{figure}[H]
  \centering
  % Levels of analysis (after McGuire, Gallagher & Ziemian, Fig. 8.1)
\begin{tikzpicture}[>=Latex,font=\footnotesize]
  \begin{axis}[
    width=13.5cm,height=9.5cm,
    axis lines=left,
    xlabel={Lateral deflection, $\Delta$},
    ylabel={Lateral load, $H$},
    xmin=0,xmax=10,ymin=0,ymax=10,
    xtick=\empty,
    ytick={8.6,7.6,5.2,3.9,2.8},
    yticklabels={$H_{ec}$,$H_{es}$,$H_{ic}$,$H_p$,$H_{is}$},
    clip=false,
  ]
    % horizontal limit lines
    \draw[dashed] (axis cs:0,8.6) -- (axis cs:9.6,8.6);
    \node[above right,inner sep=1pt] at (axis cs:6.2,8.6) {Elastic critical load};
    \draw[dashed] (axis cs:0,7.6) -- (axis cs:9.6,7.6);
    \node[above right,inner sep=1pt] at (axis cs:6.2,7.6) {Elastic stability limit};
    \draw[dashed] (axis cs:0,5.2) -- (axis cs:9.6,5.2);
    \node[above right,inner sep=1pt] at (axis cs:6.2,5.2) {Inelastic critical load};
    \draw[dashed] (axis cs:0,3.9) -- (axis cs:9.6,3.9);
    \node[above right,inner sep=1pt] at (axis cs:6.2,3.9) {Plastic limit load};
    \draw[dashed] (axis cs:0,2.8) -- (axis cs:9.6,2.8);
    \node[above right,inner sep=1pt] at (axis cs:6.2,2.8) {Inelastic stability limit};
    % first-order elastic: straight line
    \draw[thick] (axis cs:0,0) -- (axis cs:3.4,10);
    \node[right,align=left,fill=white,inner sep=1pt] at (axis cs:3.5,9.4) {First-order\\elastic analysis};
    % second-order elastic: softening, asymptote to elastic stability limit
    \draw[thick] (axis cs:0,0) .. controls (axis cs:1.4,4.2) and (axis cs:3.0,7.4) .. (axis cs:6.2,7.6)
                 .. controls (axis cs:7.0,7.6) .. (axis cs:7.6,7.3);
    % bifurcation branch at H_ec (elastic)
    \draw[thick] (axis cs:0.9,8.6) -- (axis cs:3.7,8.6);
    \node[left] at (axis cs:2.6,8.95) {Bifurcation};
    % bifurcation on 2nd-order elastic curve
    \draw[thick] (axis cs:3.9,6.6) .. controls (axis cs:4.6,6.4) .. (axis cs:5.4,6.1);
    \node[right] at (axis cs:4.0,6.95) {Bifurcation};
    \node[right,align=left,fill=white,inner sep=1pt] at (axis cs:4.3,5.4) {Second-order\\elastic analysis};
    % first-order inelastic: follows first-order elastic then flattens to H_p
    \draw[thick] (axis cs:0,0) -- (axis cs:1.05,3.1) .. controls (axis cs:1.6,3.8) .. (axis cs:2.6,3.9)
                 -- (axis cs:5.8,3.9);
    \node[right,align=left,fill=white,inner sep=1pt] at (axis cs:3.2,4.5) {First-order\\inelastic analysis};
    % bifurcation at H_ic (inelastic)
    \draw[thick] (axis cs:0.6,5.2) -- (axis cs:2.6,5.2);
    \node[left] at (axis cs:1.6,5.5) {Bifurcation};
    % second-order inelastic: peaks at H_is then descends
    \draw[thick] (axis cs:0,0) -- (axis cs:0.85,2.5) .. controls (axis cs:1.2,2.8) and (axis cs:1.8,2.9) .. (axis cs:2.6,2.8)
                 .. controls (axis cs:3.4,2.6) .. (axis cs:4.4,1.7);
    \node[right,align=left] at (axis cs:1.7,1.5) {Second-order\\inelastic analysis};
    % small inset frame
    \begin{scope}[shift={(axis cs:7.4,0.4)},x=0.22cm,y=0.22cm]
      \draw[thick] (0,0) -- (0,7) -- (5,7) -- (5,0);
      \draw[thick] (0,3.5) -- (5,3.5);
      \fill[pattern={Lines[angle=45,distance=2pt]}] (-0.6,-0.5) rectangle (0.6,0);
      \fill[pattern={Lines[angle=45,distance=2pt]}] (4.4,-0.5) rectangle (5.6,0);
      \draw[->] (-1.8,7) -- (-0.1,7) node[pos=0,left] {\tiny $H=\alpha P$};
      \draw[->] (-1.8,3.5) -- (-0.1,3.5) node[pos=0,left] {\tiny $H=\alpha P$};
      \draw[->] (0.6,8.6) -- (0.6,7.2) node[pos=0,above] {\tiny $P$};
      \draw[->] (4.4,8.6) -- (4.4,7.2);
      \draw[->] (0.6,5.1) -- (0.6,3.7) node[pos=0,above] {\tiny $P$};
      \draw[->] (4.4,5.1) -- (4.4,3.7);
      \draw[<->] (5.3,0.2) -- (5.3,7) ; \node[right] at (5.3,3.5) {\tiny $\Delta$};
    \end{scope}
  \end{axis}
\end{tikzpicture}

  \caption{Levels of analysis (after Fig.~8.1 of the textbook).}
\end{figure}
